%!TEX root = ../Thesis.tex

\chapter{Konzeptionelle Modellierung und Design}\label{cap:KonzepModelDesign}
\markboth{Konzeptionelle Modellierung und Design}{Konzeptionelle Modellierung und Design}

Kapitel 3 konkretisiert den Entwurf des prototypischen Systems als konzeptionelles Artefakt und beschreibt die dafür notwendigen Modelle in einer konsistenten, aufeinander aufbauenden Darstellung. \cite{nunamaker1990}
Methodisch wird hierfür die Phase des Theory Building nach Nunamaker Jr., Chen und Purdin herangezogen und in dieser Arbeit auf das Forschungsziel FZ2T angewendet, da in dieser Phase insbesondere Konzepte, Modelle und konzeptionelle Rahmen zur Beschreibung eines Artefakts entwickelt und präzisiert werden. \cite{nunamaker1990}
Die Struktur des Kapitels orientiert sich zudem am User Centered System Design nach Norman und Draper, indem die beabsichtigte Nutzung und die Aufgaben der Zielgruppe als Ausgangspunkt für die Modellbildung gesetzt werden. \cite{norman_draper1986}

Das Vorgehen folgt einer Ableitungslogik, bei der zunächst die Nutzungssicht über ein konzeptionelles Anwendungsfallmodell beschrieben wird, um Ziele, Interaktionen und Systemleistungen aus Sicht der primären Nutzerrolle eindeutig festzulegen. \cite{norman_draper1986}
Auf dieser Basis wird ein konzeptionelles Informationsmodell entwickelt, das die für Erhebung, Bewertung und Ergebnisbildung benötigten Informationsobjekte und Beziehungen definiert und damit die semantische Grundlage der Verarbeitung festlegt. \cite{nunamaker1990}
Abschließend wird ein konzeptionelles Architekturmodell formuliert, das die Systemverantwortlichkeiten in Komponenten strukturiert und die Kommunikation zwischen den Komponenten so beschreibt, dass die zuvor definierten Anwendungsfälle und Informationsobjekte im Entwurf nachvollziehbar realisiert werden können. \cite{nunamaker1990}

Die Voraussetzungen, Annahmen und die detaillierte Systemabgrenzung werden im anschließenden Abschnitt zum konzeptionellen Anwendungsfallmodell eingeführt, damit die folgenden Modelle auf einem einheitlichen Verständnis von Nutzungskontext und Systemgrenze aufsetzen. \cite{norman_draper1986}


\section{Konzeptionelles Anwendungsfallmodell}\label{sec:konzeptionellesAnwendungsfallmodell}
Das konzeptionelle Anwendungsfallmodell bildet den Ausgangspunkt für die nachfolgenden konzeptionellen Modelle, indem die beabsichtigte Nutzung des prototypischen Systems aus Sicht der primären Nutzerrolle beschrieben wird. \cite{jacobson2024_usecase3}
Die Modellierung orientiert sich methodisch an der Phase des Theory Building nach Nunamaker Jr., Chen und Purdin, da in dieser Phase der Entwurf des Artefakts in Form von konzeptionellen Beschreibungen und Modellen systematisch konkretisiert wird. \cite{nunamaker1990}
Die Struktur dieses Kapitels folgt darüber hinaus der Logik des User Centered System Design nach Norman und Draper, indem der Nutzen und die Ziele der Nutzer als leitendes Prinzip für die Ableitung von Systemleistungen und Modellgrenzen gesetzt werden. \cite{norman_draper1986}
In diesem Sinne wird zunächst festgelegt, welche Ziele der Nutzer mit dem System erreicht, welche Systemleistungen hierfür benötigt und welche Interaktionen dafür vorgesehen werden. \cite{norman_draper1986}
Die Ergebnisse werden als UML-Anwendungsfalldiagramm und als konsistenter Katalog von Use Cases dokumentiert, um die Systemgrenze und die funktionalen Leistungen präzise und anschlussfähig für die weiteren Modelle zu definieren. \cite{omg2017_uml251_about,jacobson2024_usecase3}


\subsection{Zielsetzung und Ausgangslage}
Ziel des konzeptionellen Anwendungsfallmodells ist es, die vom System zu erbringenden Leistungen so zu beschreiben, dass die spätere Ableitung eines Informationsmodells und eines Architekturmodells auf einer stabilen, nutzerorientierten Grundlage erfolgt. \cite{jacobson2024_usecase3}
Das Anwendungsfallmodell konkretisiert hierzu die beabsichtigten Nutzungsszenarien, indem es die Nutzerziele, die Auslöser der Interaktion sowie die erwarteten Ergebnisse des Systems in einer strukturierten Form festhält. \cite{jacobson2024_usecase3}
Als methodischer Rahmen wird die Modellbildung als Bestandteil des Theory Building verstanden, sodass das Anwendungsfallmodell nicht nur dokumentiert, sondern als bewusste Artefaktkonstruktion im Forschungsprozess begründet wird. \cite{nunamaker1990}
Die Ausgangslage wird dadurch bestimmt, dass KMU für die kombinierte Einführung von \gls{bi} und Prozessautomatisierung eine strukturierte Entscheidungs- und Priorisierungsunterstützung benötigen, die auf einem systematisierten Vorgehen beruht und über einzelne Toolentscheidungen hinausgeht. \cite{norman_draper1986}
Das zu konzipierende System wird in diesem Kapitel als prototypisches, eigenständiges System modelliert, das die Erhebung eines unternehmensspezifischen Use Cases und die Bearbeitung eines strukturierten Fragenkatalogs als primären Interaktionsmodus bereitstellt. \cite{jacobson2024_usecase3}
Die Verarbeitung der erhobenen Informationen ist so ausgerichtet, dass zwei getrennte Bewertungen erzeugt werden, die eine \gls{bi}-bezogene Reifegradbewertung und eine Reifegradbewertung der Prozessautomatisierung, deren Ergebnisse anschließend zu einer integrierten Sicht zusammengeführt werden. \cite{nunamaker1990}
Der zentrale Output des Systems wird als priorisierter Maßnahmenkatalog verstanden, der aus den Bewertungen und deren Synthese abgeleitet wird und dem Nutzer konkrete Handlungsoptionen für die Weiterentwicklung liefert. \cite{jacobson2024_usecase3}
Die Rolle eines \gls{llm} wird im System als assistierende Unterstützung für die Interpretation der Eingaben und die Formulierung von Vorschlägen konzipiert, während die finale Entscheidung über Auswahl und Priorisierung der Maßnahmen beim Nutzer verbleibt. \cite{norman_draper1986}
Die Systemgrenze wird im Anwendungsfallmodell bewusst als geschlossen definiert, sodass keine direkten Schnittstellen zu operativen Unternehmenssystemen vorgesehen werden und sämtliche Informationen über die Nutzerinteraktion eingebracht werden. \cite{jacobson2024_usecase3}
Als primäre Nutzerrolle wird ein „KMU-Entscheider beziehungsweise Architekt“ angesetzt, der die Planung und Priorisierung von BI- und Automatisierungsvorhaben verantwortet und das System zur strukturierten Vorbereitung strategischer Entscheidungen nutzt. \cite{norman_draper1986}
Die grundlegenden Annahmen zur Systemnutzung und zur Abgrenzung werden in diesem Abschnitt verankert, während die detaillierte Ausprägung der Vorbedingungen und Alternativabläufe in den einzelnen Use-Case-Beschreibungen konkretisiert wird. \cite{jacobson2024_usecase3}
Die Wahl von UML als Notation stellt sicher, dass die im Use-Case-Modell festgehaltenen Interaktionen als standardisierte Darstellung in die weiteren konzeptionellen Sichten überführt werden können. \cite{omg2017_uml251_about}


\subsection{Aufgabenanalyse und Ableitung der Systemleistungen}
Die Aufgabenanalyse folgt dem Grundgedanken des User Centered System Design, wonach die Gestaltung eines Systems von den Zielen, Aufgaben und Bedürfnissen der Nutzer ausgeht und erst nachgelagert in konkrete Systemleistungen überführt wird. \cite{normanDraper1986}
Für die Strukturierung der Aufgaben wird der Nutzungskontext explizit als Ausgangspunkt gesetzt, weil sich Anforderungen an Interaktion und Ergebnisdarstellung aus dem Arbeitsumfeld und den Entscheidungsprozessen der Zielgruppe ableiten. \cite{normanDraper1986}
Dieses Vorgehen entspricht zugleich einem human-centred Designverständnis, in dem Aufgabenorientierung und Kontextbezug als zentrale Aktivitäten der Gestaltung verankert sind. \cite{iso9241_210_2019}

Aus der Perspektive der primären Nutzerrolle wird das übergeordnete Ziel als strategische Vorbereitung und Priorisierung von Vorhaben zur kombinierten Einführung von \gls{bi} und Prozessautomatisierung verstanden.
Dieses Ziel wird in eine sequenzielle Aufgabenlogik zerlegt, um die Systemunterstützung entlang eines nachvollziehbaren Arbeitsablaufs zu bestimmen und damit die Grundlage für das spätere UML-Use-Case-Modell zu schaffen. \cite{normanDraper1986}
Die Zerlegung orientiert sich an der Idee, dass Systeme nicht einzelne Funktionen isoliert bereitstellen, sondern vollständige, für den Nutzer nützliche Unterstützungseinheiten entlang seiner Aufgaben liefern sollen. \cite{iso9241_210_2019}

Die erste Aufgabenphase umfasst die strukturierte Erhebung des betrachteten Use Cases und des dazugehörigen Kontextwissens, weil die spätere Bewertung eine hinreichend präzise Problem- und Situationsbeschreibung voraussetzt. \cite{iso9241_210_2019}
Daraus ergibt sich als Systemleistung die geführte Erfassung des Use Cases sowie die Durchführung eines standardisierten Fragenkatalogs, der die relevanten Merkmale in einer vergleichbaren Form erhebt.
Damit die Ergebnisse belastbar weiterverarbeitet werden können, umfasst diese Leistung eine formale Plausibilisierung und Vollständigkeitsprüfung der Eingaben, weil inkonsistente oder unvollständige Antworten die nachgelagerten Bewertungen systematisch verzerren können. \cite{iso9241_210_2019}
Als Ergebnisartefakt der Erhebungsphase entsteht ein konsolidierter Satz an Antworten, der die Grundlage für beide Reifegradbewertungen bildet.

Die zweite Aufgabenphase adressiert die Einordnung des Status quo, weil der Nutzer eine systematische Ausgangsbasis benötigt, um Handlungsbedarfe nicht ausschließlich intuitiv abzuleiten. \cite{normanDraper1986}
Da in dieser Arbeit zwei Reifegradmodelle vorgesehen sind, wird die Bewertung in eine \gls{bi}-bezogene Teilbewertung und eine Teilbewertung für Prozessautomatisierung differenziert, um domänenspezifische Stärken und Defizite getrennt sichtbar zu machen.
Als Systemleistung wird daher benötigt, die erhobenen Eingaben jeweils in eine nachvollziehbare Reifegradeinschätzung zu überführen und die Ergebnisse in einer Form zurückzugeben, die der Nutzer unmittelbar in strategische Entscheidungen überführen kann. \cite{normanDraper1986}
Die Verständlichkeit der Bewertungen ist dabei eine zentrale Anforderung, weil UCD die Nutzbarkeit eines Ergebnisses an dessen Passung zu Zielen und Aufgaben der Nutzer festmacht. \cite{iso9241_210_2019}

Die dritte Aufgabenphase bündelt beide Teilbewertungen in einer Synthese, weil die kombinierte Einführung von \gls{bi} und Prozessautomatisierung Abhängigkeiten erzeugt, die in getrennten Bewertungen nur unzureichend sichtbar werden.
Die dazugehörige Systemleistung besteht in der konsistenten Zusammenführung der beiden Bewertungen sowie der Explikation relevanter Wechselwirkungen, sodass Prioritäten und Voraussetzungen für Maßnahmen aus einer integrierten Sicht abgeleitet werden können. \cite{normanDraper1986}
Eine solche Synthese entspricht dem Anspruch, im Theory Building konzeptionelle Strukturen zu entwickeln, die den Zusammenhang zwischen Problem, Artefaktlogik und Ergebnissen nachvollziehbar beschreiben. \cite{nunamaker1990}

Die vierte Aufgabenphase umfasst die Ableitung eines priorisierbaren Maßnahmenkatalogs, weil die Nutzerrolle ein operationalisierbares Ergebnis benötigt, das als Grundlage für Planung, Roadmapping und interne Abstimmung genutzt werden kann.
Daraus folgt als Systemleistung die Generierung eines initialen Maßnahmenvorschlags, der sich aus den beiden Reifegradbewertungen und der Synthese begründet ableitet und als handlungsorientiertes Ergebnis aufbereitet ist.
Da das System assistierend konzipiert ist, muss eine zusätzliche Systemleistung vorgesehen werden, die dem Nutzer eine Anpassung und Priorisierung erlaubt, um Kontextwissen, Restriktionen und Präferenzen in die finale Auswahl zu übertragen. \cite{iso9241_210_2019}
Diese Nutzerkontrolle ist konsistent mit dem UCD-Prinzip, dass die Gestaltung Interaktion so unterstützen soll, dass der Nutzer Entscheidungen informiert und selbstbestimmt treffen kann. \cite{normanDraper1986}

Über alle Aufgabenphasen hinweg wird das \gls{llm} als internes Unterstützungselement verstanden, das die semantische Interpretation der Antworten und die Formulierung von Vorschlägen unterstützt, ohne die Ergebnisverantwortung der Nutzerrolle zu ersetzen.
Damit wird die Aufgabenanalyse zugleich zur Klammer zwischen nutzerzentrierter Herleitung und formaler Modellierung, weil sie die benötigten Systemleistungen als Ergebnis einer expliziten Nutzer- und Aufgabenperspektive begründet. \cite{iso9241_210_2019}

\subsection{UML-Use-Case-Diagramm als formale Konsolidierung}
Das UML-Use-Case-Diagramm konsolidiert die in der Aufgabenanalyse identifizierten Nutzeraufgaben zu expliziten, wertorientierten Systemleistungen und übersetzt diese in eine standardisierte Notation, die unabhängig von konkreten Implementierungsentscheidungen interpretierbar bleibt. \cite{jacobson_usecase3_2024}
Die UML dient in diesem Zusammenhang als formale Modellierungssprache zur Visualisierung und Dokumentation von Systemaspekten und ermöglicht eine konsistente Überführung der Nutzungssicht in anschlussfähige Folgemodelle wie Informations- und Architekturmodelle. \cite{omg_uml251_about}
Die Strukturierung entlang eines nutzerzentrierten Vorgehens bleibt dabei erhalten, indem Use Cases nicht als technische Funktionen, sondern als vom Nutzer erwartete, beobachtbare Ergebnisse des Systems formuliert werden. \cite{iso9241_210_2019_sample,jacobson_usecase3_2024}

\begin{figure}[H]
\centering
\includegraphics[width=14cm]{arbeit/images/UML-UseCase-Diagramm.png}
\caption[Use-Case-Diagramm]{UML-Use-Case-Diagramm zur Darstellung der Nutzerinteraktionen und Systemleistungen}
\label{IMG:UML-UseCase-Diagramm}
\end{figure}

Das dargestellte Diagramm fasst den End-to-End-Nutzen in einem übergeordneten Anwendungsfall zusammen, der die Ableitung von Bewertungen und Maßnahmen als zentrales Nutzerziel abbildet. \cite{jacobson_usecase3_2024}
Die Systemgrenze wird als geschlossenes prototypisches System modelliert, sodass ausschließlich die Interaktion zwischen der primären Nutzerrolle und dem System dargestellt wird und keine externen Unternehmenssysteme als Akteure auftreten. \cite{uml_diagrams_usecase_overview}
Die Single-Role-Entscheidung wird im Diagramm durch einen einzigen Akteur „KMU-Entscheider/Architekt“ umgesetzt, der die systematische Planung und Priorisierung von \gls{bi}- und Automatisierungsvorhaben repräsentiert. \cite{jacobson_usecase3_2024}
Das \gls{llm} wird nicht als Akteur modelliert, weil es als internes Verarbeitungselement innerhalb der Systemgrenze konzipiert ist und die Diagrammart primär externe Interaktionsbeziehungen zwischen Nutzern und System fokussiert. \cite{uml_diagrams_usecase_overview}

Die Beziehungen zwischen den Use Cases werden so gewählt, dass Pflichtbestandteile des Workflows als \texttt{<<include>>}-Beziehungen modelliert sind, um wiederverwendbare und für das Erreichen des Nutzerziels notwendige Teilschritte explizit auszuweisen. \cite{uml_diagrams_include,visual_paradigm_uc_notations}
Dies betrifft insbesondere die Erhebungs- und Bewertungsleistungen, da der Maßnahmenkatalog in der gewählten Systemausprägung zwingend auf dem erfassten Use Case, dem beantworteten Fragenkatalog sowie den beiden Teilbewertungen für \gls{bi} und Prozessautomatisierung basiert. \cite{jacobson_usecase3_2024}
Optionale, vom Nutzer situativ angestoßene Schritte werden als \texttt{<<extend>>}-Beziehungen modelliert, um den assistierenden Charakter des Systems abzubilden, bei dem Vorschläge generiert werden, die finale Entscheidung und Priorisierung jedoch beim Nutzer verbleibt. \cite{visual_paradigm_uc_notations}
Auf dieser Basis wird die Anpassung und Priorisierung als Erweiterung der initialen Generierung des Maßnahmenkatalogs verstanden, während ein Export als nachgelagerte, ebenfalls optionale Nutzung des Ergebnisses modelliert wird. \cite{cockburn_effective_use_cases_pdf}
Die im Diagramm ergänzte Notiz zur assistierenden Verarbeitung durch das \gls{llm} fungiert als konzeptionelle Klarstellung der Verantwortungsverteilung, ohne die Rollenlogik des Use-Case-Modells zu erweitern, und unterstützt damit die spätere konsistente Ableitung von Informationsobjekten und Komponentenverantwortlichkeiten. \cite{iso9241_210_2019_sample}


\subsection{Auswahl der genutzten Reifegrad- und Readiness-Modelle}
\label{subsec:kap3_auswahl_maturity_readiness}

Die Auswahl der in dieser Arbeit verwendeten Reifegrad- und Readiness-Modelle folgt einem transparent dokumentierten Such- und Bewertungsprozess, um Nachvollziehbarkeit und Wiederverwendbarkeit der konzeptionellen Modellierung sicherzustellen.\cite{vomBrocke2009Reconstructing,WebsterWatson2002LiteratureReview,OkoliSchabram2010SLRGuide}
Die Literatursuche wurde als strukturierte Recherche konzipiert, da entsprechende Leitfäden die explizite Dokumentation von Suchstrings, Datenquellen und Selektionsschritten als Voraussetzung für die Beurteilung der Vollständigkeit einer Recherche benennen.\cite{vomBrocke2009Reconstructing,OkoliSchabram2010SLRGuide,KitchenhamCharters2007SLRGuidelines}
Für die Operationalisierung der Auswahlkriterien wurden etablierte Vorgehensmodelle zur Entwicklung und Bewertung von Reifegradmodellen herangezogen, da diese explizit Anforderungen an Zweckbezug, Dimensionen, Messbarkeit und Stufenlogik formulieren.\cite{deBruin2005MainPhases,BeckerKnackstedtPoeppelbuss2009ProcedureModel,PoeppeelbussRoeglinger2011DesignPrinciples}

Als Selektionskriterien für die Modellbasis wurden erstens die inhaltliche Passung zu den beiden Untersuchungsbereichen Data \& Analytics beziehungsweise Prozessautomatisierung, zweitens die Operationalisierbarkeit in Items für einen Fragebogen, drittens die Abdeckung sowohl technischer als auch organisatorischer Fähigkeiten sowie viertens die Anschlussfähigkeit an die in Kapitel~2 abgeleiteten Gestaltungskomponenten definiert.\cite{BeckerKnackstedtPoeppelbuss2009ProcedureModel,PoeppeelbussRoeglinger2011DesignPrinciples,deBruin2005MainPhases}
Diese Kriterien entsprechen der Zielsetzung, ein Messinstrument zu entwerfen, das Dimensionen konsistent aggregiert und zugleich konkrete Verbesserungshinweise ermöglicht.\cite{PoeppeelbussRoeglinger2011DesignPrinciples,BeckerKnackstedtPoeppelbuss2009ProcedureModel}
Zur Identifikation geeigneter Kandidaten wurden Übersichtsarbeiten zu Data-\&-Analytics-Reifegradmodellen sowie zu \gls{rpa}-Adoption und organisatorischen Erfolgsfaktoren priorisiert, da Reviews eine verdichtete Sicht auf wiederkehrende Domänen und Konstrukte liefern und damit die Grundlage für eine belastbare Dimensionsableitung bilden.\cite{Langer2025DataAnalyticsMaturityReview,ChuahWong2011BIMaturityReview,daSilvaCostaMamedeSilva2022RPAAdoptionSLR}
Für die Auswahlentscheidung wurden anschließend die in den Verfahren zur Reifegradmodellkonstruktion beschriebenen Artefaktmerkmale genutzt, insbesondere Klarheit des Anwendungszwecks, explizite Dimensionendefinition, Vorliegen einer Assessment-Logik sowie die Möglichkeit, Dimensionen in einzelne Items zu überführen.\cite{deBruin2005MainPhases,BeckerKnackstedtPoeppelbuss2009ProcedureModel,PoeppeelbussRoeglinger2011DesignPrinciples}

Für den Bereich Data \& Analytics wurden als Modellbasis primär die in der systematischen Übersichtsarbeit zu Data-\&-Analytics-Maturity-Modellen zusammengefassten Domänen und Strukturmerkmale verwendet, da diese Arbeit eine konsolidierte Sicht auf wiederkehrende Reifegraddomänen liefert und mehrere Dutzend Modelle vergleichend auswertet.\cite{Langer2025DataAnalyticsMaturityReview}
Ergänzend wurde ein domänenspezifisches \gls{bia}-Reifegradmodell aus einer Design-Science-Studie herangezogen, da dieses eine explizite Assessment-Methode und operationalisierbare Domänen beschreibt und damit die Itemableitung für den Fragebogen unterstützt.\cite{CardosoSu2022BIAMaturityModel,BeckerKnackstedtPoeppelbuss2009ProcedureModel}
Die Kombination aus systematischer Review-Perspektive und konkret ausformulierter Assessment-Methode reduziert das Risiko, ausschließlich vendorgetriebene oder unvollständig dokumentierte Reifegradansätze zu übernehmen, das in Übersichtsarbeiten zu BI-Maturity-Modellen als Problem beschrieben wird.\cite{ChuahWong2011BIMaturityReview,Langer2025DataAnalyticsMaturityReview}

Für den Bereich Prozessautomatisierung wurde die Modellbasis aus einer systematischen Literaturübersicht zur \gls{rpa}-Adoption abgeleitet, da diese wiederkehrende Erfolgsfaktoren, Barrieren und Eignungskriterien für Prozesse syntheseartig zusammenführt und damit direkt in ein Readiness-Assessment überführbar ist.\cite{daSilvaCostaMamedeSilva2022RPAAdoptionSLR,KitchenhamCharters2007SLRGuidelines}
Zur Präzisierung organisatorischer Anforderungen an Skalierung und nachhaltigen Betrieb wurde zusätzlich Literatur zu Governance in \gls{rpa}-Programmen herangezogen, da Governance-Strukturen als wiederkehrender Gestaltungs- und Erfolgsfaktor bei \gls{rpa} beschrieben werden.\cite{KedzioraPenttinen2021RPAGovernance,daSilvaCostaMamedeSilva2022RPAAdoptionSLR}
Als übergreifender Bezugsrahmen für Readiness-Konstrukte wurde das Technology-Organization-Environment-Framework genutzt, da es organisationale Adoption entlang technologischer, organisationaler und umweltbezogener Einflussgrößen strukturiert und in firm-level Adoption Reviews als Referenzrahmen berichtet wird.\cite{TornatzkyFleischer1990TOE,OliveiraMartins2011ITAdoptionReview}
Damit basiert die in Abschnitt~\ref{subsec:kap3_auswahl_maturity_readiness} definierte Modellbasis auf einer Kombination aus Reifegrad- und Readiness-Perspektive, die sowohl Fähigkeitsausprägungen als auch Einführungsvoraussetzungen abbildet.\cite{PoeppeelbussRoeglinger2011DesignPrinciples,TornatzkyFleischer1990TOE,Langer2025DataAnalyticsMaturityReview}
Die konkrete Operationalisierung dieser Modellbasis in Dimensionen, Items sowie Scoring- und Output-Logik wird in den nachfolgenden Abschnitten des Kapitels als Bestandteil des konzeptionellen Modells spezifiziert und in Kapitel~4 prototypisch umgesetzt.\cite{deBruin2005MainPhases,BeckerKnackstedtPoeppelbuss2009ProcedureModel}


\subsection{Herleitung und Operationalisierung des Fragenkatalogs}
\label{subsec:kap3_herleitung_fragenkatalog}

Die Herleitung des Fragenkatalogs folgt einem deduktiven Vorgehen, bei dem Items aus theoretisch definierten Inhaltsdomänen abgeleitet werden, um die in Abschnitt~\ref{subsec:kap3_auswahl_maturity_readiness} festgelegten Reifegrad- und Readiness-Konstrukte in beobachtbare Indikatoren zu überführen.\cite{Hinkin1998ScaleDevelopment,deBruin2005MainPhases,BeckerKnackstedtPoeppelbuss2009ProcedureModel}
Ein deduktives Item-Design setzt eine explizite Definition der Konstrukte und ihres Geltungsbereichs voraus, da die spätere Interpretation der Messwerte an die Abdeckung der Inhaltsdomäne gebunden ist.\cite{Hinkin1998ScaleDevelopment,Boateng2018ScaleDevelopment}
Die in Kapitel~2 eingeführten Maturity- und Readiness-Perspektiven werden hierfür als theoretische Grundlage genutzt, weil Reifegradmodelle Dimensionen zur Fähigkeitsausprägung strukturieren und Readiness-Rahmenwerke Voraussetzungen der Einführung und Skalierung systematisieren.\cite{PoeppeelbussRoeglinger2011DesignPrinciples,TornatzkyFleischer1990TOE,OliveiraMartins2011ITAdoptionReview}
Für Data \& Analytics werden die in Reviews identifizierten wiederkehrenden Domänen als Itemdomänen interpretiert, da Reviews Domänen über mehrere Modelle hinweg konsolidieren und damit eine stabile Grundlage für die Dimensionsbildung liefern.\cite{Langer2025DataAnalyticsMaturityReview,CardosoSu2022BIAMaturityModel}
Für Prozessautomatisierung werden Faktoren zur \gls{rpa}-Adoption und \gls{rpa}-Governance als Itemdomänen genutzt, da diese Faktoren in systematischen Übersichten und Fallstudien als wiederkehrende Barrieren, Erfolgsfaktoren und Skalierungsvoraussetzungen beschrieben werden.\cite{daSilvaCostaMamedeSilva2022RPAAdoptionSLR,KedzioraPenttinen2021RPAGovernance}

Die Operationalisierung erfolgt in drei Schritten, um eine nachvollziehbare Traceability von Framework-Domänen zu Fragebogenitems sicherzustellen.\cite{Hinkin1998ScaleDevelopment,Boateng2018ScaleDevelopment}
Im ersten Schritt werden aus den ausgewählten Modellbasen Dimensionen gebildet, die jeweils eine klar abgegrenzte Fähigkeit oder Voraussetzung repräsentieren.\cite{deBruin2005MainPhases,BeckerKnackstedtPoeppelbuss2009ProcedureModel}
Im zweiten Schritt werden pro Dimension mehrere Itemkandidaten formuliert, um die Inhaltsdomäne breit zu sampeln und eine spätere Reduktion ohne Verlust zentraler Inhalte zu ermöglichen.\cite{Hinkin1998ScaleDevelopment,Boateng2018ScaleDevelopment}
Im dritten Schritt werden die Items auf Eindeutigkeit, Messbarkeit und Praxisbezug überprüft, da diese Eigenschaften als Voraussetzungen für die Interpretierbarkeit von Assessments beschrieben werden.\cite{PoeppeelbussRoeglinger2011DesignPrinciples,Boateng2018ScaleDevelopment}

Zur Absicherung der Inhaltsvalidität wird eine Expertenbeurteilung vorgesehen, da Content-Validation-Ansätze die Eignung einzelner Items zur Abdeckung einer Inhaltsdomäne durch strukturierte Expertenurteile bewerten.\cite{Lawshe1975ContentValidity,AyreScally2014CVR,Boateng2018ScaleDevelopment}
Als quantifizierende Ergänzung wird die Content Validity Ratio nach \citeauthor{Lawshe1975ContentValidity} genutzt, weil die CVR die Essentialität eines Items anhand des Anteils zustimmender Expert:innen operationalisiert und damit eine transparente Entscheidungsregel für Revision oder Ausschluss einzelner Items bereitstellt.\cite{Lawshe1975ContentValidity,AyreScally2014CVR}
Die resultierende Itemmenge wird anschließend als versioniertes Artefakt geführt, da Reifegradassessments eine stabile und nachvollziehbar evolvierende Instrumentenbasis benötigen, um Ergebnisse über Zeitpunkte und Kontexte hinweg vergleichbar zu halten.\cite{Stoiber2023KeepingMaturityAlive,deBruin2005MainPhases}



\subsection{Scoring- und Threshold-Modell sowie Maßnahmenlogik}
\label{subsec:kap3_scoring_threshold_massnahmen}

Das Scoring-Modell operationalisiert die im Fragenkatalog definierten Items als messbare Dimensionsergebnisse und stellt damit die Grundlage für die in Abschnitt~\ref{subsec:kap3_output_artefakte} spezifizierten Assessment-Outputs dar.\cite{deBruin2005MainPhases,BeckerKnackstedtPoeppelbuss2009ProcedureModel}
Eine dimensionsbasierte Aggregation wird genutzt, weil Reifegradmodelle Fähigkeiten typischerweise über Domänen beziehungsweise Dimensionen strukturieren und Assessments dadurch interpretierbarer werden.\cite{deBruin2005MainPhases,PoeppeelbussRoeglinger2011DesignPrinciples}
Die Spezifikation des Scoring-Modells in Kapitel~3 reduziert Implementierungsfreiheit in Kapitel~4 und unterstützt die Nachvollziehbarkeit der Artefaktkonstruktion im Sinne der Design-Science-Leitlinien.\cite{Hevner2004DesignScience,MarchSmith1995DesignNaturalScience}

\paragraph{Dimensionen und Item-Zuordnung.}
Jedes Item des Fragenkatalogs wird genau einer Bewertungsdimension zugeordnet, da eine eindeutige Zuordnung die spätere Aggregation und Interpretation stabilisiert.\cite{deBruin2005MainPhases,BeckerKnackstedtPoeppelbuss2009ProcedureModel}
Die Dimensionen werden entlang der zuvor ausgewählten Reifegrad- und Readiness-Modellbasis gebildet, um inhaltliche Konsistenz zwischen Theoriebezug und Messinstrument herzustellen.\cite{PoeppeelbussRoeglinger2011DesignPrinciples,TornatzkyFleischer1990TOE,OliveiraMartins2011ITAdoptionReview}
Der Fragenkatalog liefert die Itemmenge $I_d$ je Dimension $d$ und ist damit die primäre Spezifikation für Scoring und Schwellenwerte.\cite{deBruin2005MainPhases,BeckerKnackstedtPoeppelbuss2009ProcedureModel}

\paragraph{Berechnung der Dimensionsscores.}
Für ein Item $i$ wird die Antwort $r_i$ auf einer ordinalen Skala erfasst und für die Aggregation in einen normierten Wert $s_i$ überführt.\cite{BeckerKnackstedtPoeppelbuss2009ProcedureModel,deBruin2005MainPhases}
Bei einer 5-stufigen Skala $r_i \in \{1,2,3,4,5\}$ wird eine lineare Normierung verwendet, um Dimensionsscores vergleichbar zu machen.\cite{BeckerKnackstedtPoeppelbuss2009ProcedureModel}
\begin{equation}
s_i = \frac{r_i - 1}{4}
\end{equation}
Der Dimensionsscore $S_d$ wird als arithmetisches Mittel der normierten Itemwerte berechnet und auf einen Bereich von $0$ bis $100$ skaliert, um eine einheitliche Interpretation über Dimensionen hinweg zu ermöglichen.\cite{deBruin2005MainPhases,BeckerKnackstedtPoeppelbuss2009ProcedureModel}
\begin{equation}
S_d = \left(\frac{1}{|I_d|}\sum_{i \in I_d} s_i \right)\cdot 100
\end{equation}
Für den PoC wird eine Gleichgewichtung der Items innerhalb einer Dimension verwendet, weil in frühen Artefaktversionen eine einfache, transparente Aggregation die Nachvollziehbarkeit gegenüber komplexeren Gewichtungsmodellen priorisiert.\cite{PoeppeelbussRoeglinger2011DesignPrinciples,BeckerKnackstedtPoeppelbuss2009ProcedureModel}
Eine spätere Erweiterung um Gewichte bleibt möglich, sofern diese Gewichte konzeptionell begründet und empirisch kalibriert werden.\cite{deBruin2005MainPhases,PoeppeelbussRoeglinger2011DesignPrinciples}

\paragraph{Thresholds und Reifegradstatus.}
Aus Dimensionsscores wird ein diskreter Reifegradstatus abgeleitet, da Reifegradmodelle Fähigkeiten häufig in Stufen beschreiben und Stufen die Kommunikation von Handlungsbedarf erleichtern.\cite{RoeglingerPoeppelbuss2011BPMMaturityModels,PoeppeelbussRoeglinger2011DesignPrinciples}
Für den Prototyp wird ein vierstufiges Schema verwendet, um eine ausreichend differenzierte, aber noch interpretierbare Abstufung zu erreichen.\cite{RoeglingerPoeppelbuss2011BPMMaturityModels}
Die Schwellenwerte werden als feste Intervalle im Bereich $0$ bis $100$ definiert, um eine eindeutige Zuordnung zu gewährleisten und die Implementierung reproduzierbar zu halten.\cite{BeckerKnackstedtPoeppelbuss2009ProcedureModel,deBruin2005MainPhases}
\begin{equation}
\text{Level}(S_d)=
\begin{cases}
L1 & \text{wenn } 0 \le S_d < 25\\
L2 & \text{wenn } 25 \le S_d < 50\\
L3 & \text{wenn } 50 \le S_d < 75\\
L4 & \text{wenn } 75 \le S_d \le 100
\end{cases}
\end{equation}
Die Verwendung fester Schwellenwerte stellt eine bewusste PoC-Entscheidung dar, da die Modellgüte in Kapitel~5 evaluiert und Schwellenwerte anschließend angepasst werden können.\cite{Hevner2004DesignScience,Nunamaker1990SystemsDevelopment}

\paragraph{Maßnahmenlogik als Mapping auf Maßnahmenklassen.}
Der Maßnahmenkatalog wird als regelbasiertes Mapping von Dimensionen und Reifegradstufen auf Maßnahmenklassen definiert, weil Reifegradmodelle in der Anwendung typischerweise zur Strukturierung von Verbesserungsinitiativen genutzt werden.\cite{RoeglingerPoeppelbuss2011BPMMaturityModels,PoeppeelbussRoeglinger2011DesignPrinciples}
Das Mapping nutzt den Reifegradstatus je Dimension als primären Trigger, da niedrige Reifegrade einen stärkeren Aufbau- und Standardisierungsbedarf signalisieren als hohe Reifegrade.\cite{RoeglingerPoeppelbuss2011BPMMaturityModels}
Die Maßnahmenklassen sind dabei als wiederverwendbare Kategorien formuliert, um Empfehlungen konsistent zu gruppieren und unabhängig von konkreten Tools beschreiben zu können.\cite{PoeppeelbussRoeglinger2011DesignPrinciples,MarchSmith1995DesignNaturalScience}

\paragraph{Definition der Maßnahmenklassen.}
Für jede Dimension $d$ wird abhängig von $\text{Level}(S_d)$ genau eine primäre Maßnahmenklasse zugeordnet, um die Outputableitung deterministisch zu halten.\cite{BeckerKnackstedtPoeppelbuss2009ProcedureModel,Hevner2004DesignScience}
\begin{itemize}
\item $L1$: \emph{Fundament schaffen} (Grundlagen, Rollen, Mindeststandards).\cite{RoeglingerPoeppelbuss2011BPMMaturityModels}
\item $L2$: \emph{Standardisieren und stabilisieren} (Prozesse, Richtlinien, Qualitätssicherung).\cite{RoeglingerPoeppelbuss2011BPMMaturityModels}
\item $L3$: \emph{Integrieren und skalieren} (Wiederverwendung, Governance, Betriebsfähigkeit).\cite{RoeglingerPoeppelbuss2011BPMMaturityModels}
\item $L4$: \emph{Optimieren} (Monitoring, kontinuierliche Verbesserung, Wertsteuerung).\cite{RoeglingerPoeppelbuss2011BPMMaturityModels}
\end{itemize}
Die konkrete inhaltliche Ausprägung einer Maßnahme wird zusätzlich an die Dimension gekoppelt, da Dimensionen die fachliche Ursache des Handlungsbedarfs repräsentieren.\cite{deBruin2005MainPhases,PoeppeelbussRoeglinger2011DesignPrinciples}

\paragraph{Priorisierung unter Berücksichtigung von Kontextrestriktionen.}
Die Priorisierung der Maßnahmen berücksichtigt neben dem Reifegradgap auch Kontextrestriktionen und Zielsetzungen aus der strukturierten Synthese, da Readiness-Ansätze Kontextfaktoren als relevante Einflussgrößen für Adoption beschreiben.\cite{TornatzkyFleischer1990TOE,OliveiraMartins2011ITAdoptionReview}
Für den PoC wird der Prioritätswert $P_d$ als Funktion aus Reifegradlücke und Kontextpassung definiert, um eine nachvollziehbare und implementierbare Regel abzubilden.\cite{BeckerKnackstedtPoeppelbuss2009ProcedureModel,Hevner2004DesignScience}
\begin{equation}
P_d = (100 - S_d)\cdot c_d
\end{equation}
Der Faktor $c_d \in [0.5,1.5]$ moduliert die Priorität abhängig von Zielkategorie und Restriktionen, damit beispielsweise stark budgetlimitierte Situationen zunächst Maßnahmenklassen mit geringerem Umsetzungsaufwand bevorzugen können.\cite{TornatzkyFleischer1990TOE,OliveiraMartins2011ITAdoptionReview}
Die konkrete Parametrisierung von $c_d$ wird als Konfigurationsartefakt umgesetzt, damit Anpassungen ohne Veränderung der Kernlogik möglich bleiben.\cite{MarchSmith1995DesignNaturalScience,Hevner2004DesignScience}

\paragraph{Ableitung der drei Output-Klassen.}
Das Ergebnisobjekt der ersten Output-Klasse enthält pro Dimension den Score $S_d$ sowie den Reifegradstatus $\text{Level}(S_d)$.\cite{deBruin2005MainPhases,BeckerKnackstedtPoeppelbuss2009ProcedureModel}
Die zweite Output-Klasse persistiert die strukturierten Kontextrestriktionen und Zielsetzungen und wird zur Modulation der Priorisierung genutzt.\cite{TornatzkyFleischer1990TOE,OliveiraMartins2011ITAdoptionReview}
Die dritte Output-Klasse erzeugt den Maßnahmenkatalog als geordnete Menge von Maßnahmenklassen je Dimension einschließlich Abhängigkeiten und Voraussetzungen, da dies die Steuerung von Verbesserungsinitiativen unterstützt.\cite{PoeppeelbussRoeglinger2011DesignPrinciples,RoeglingerPoeppelbuss2011BPMMaturityModels}




\subsection{Gestaltung der Output-Artefakte und Ableitungslogik}
\label{subsec:kap3_output_artefakte}

Die Output-Artefakte werden als Bestandteil des konzeptionellen Modells spezifiziert, da Artefaktdefinitionen in Design-Science-Arbeiten die Grundlage für Implementierung und Evaluation bilden.\cite{Hevner2004DesignScience,MarchSmith1995DesignNaturalScience,Nunamaker1990SystemsDevelopment}
Reifegradmodelle werden in der Literatur nicht nur zur Standortbestimmung verwendet, sondern auch zur Ableitung und Priorisierung von Verbesserungsmaßnahmen, wodurch Output-Artefakte über reine Score-Werte hinaus erforderlich werden.\cite{PoeppeelbussRoeglinger2011DesignPrinciples,RoeglingerPoeppelbuss2011BPMMaturityModels}
Für den vorliegenden Use Case werden daher drei Output-Klassen spezifiziert, die jeweils unterschiedliche Zwecke abdecken und gemeinsam die Nachvollziehbarkeit der Empfehlung erhöhen.\cite{Hevner2004DesignScience,PoeppeelbussRoeglinger2011DesignPrinciples}

Als erste Output-Klasse werden Assessment-Ergebnisse definiert, die pro Dimension einen Score und einen Reifegradstatus ausweisen, da eine dimensionsbasierte Aggregation die Interpretierbarkeit von Reifegradmessungen unterstützt.\cite{deBruin2005MainPhases,BeckerKnackstedtPoeppelbuss2009ProcedureModel}
Als zweite Output-Klasse wird eine strukturierte Synthese vorgesehen, die Kontextrestriktionen und Zielsetzungen explizit dokumentiert, da Readiness-Ansätze technologische, organisatorische und umweltbezogene Einflüsse als relevante Kontextbedingungen für Adoption beschreiben.\cite{TornatzkyFleischer1990TOE,OliveiraMartins2011ITAdoptionReview}
Als dritte Output-Klasse wird ein Maßnahmenkatalog spezifiziert, der Empfehlungen auf Dimensionen zurückführt, Prioritäten ausweist und Voraussetzungen sowie Abhängigkeiten benennt, da Maturity-Model-Anwendungen typischerweise Verbesserungsinitiativen strukturieren und deren Fortschritt steuerbar machen.\cite{PoeppeelbussRoeglinger2011DesignPrinciples,RoeglingerPoeppelbuss2011BPMMaturityModels}

Die Ableitungslogik wird als Regelwerk konzipiert, das von Itemantworten über Dimensionsaggregation zu Maßnahmenklassen führt, weil die explizite Spezifikation von Regeln die Reproduzierbarkeit eines Assessments erhöht.\cite{Hevner2004DesignScience,MarchSmith1995DesignNaturalScience}
Die Zuordnung von Maßnahmen zu Dimensionen wird aus der Modellbasis abgeleitet, da die verwendeten Maturity- und Readiness-Domänen die inhaltliche Struktur liefern, entlang derer Verbesserungshebel konsistent gruppiert werden können.\cite{Langer2025DataAnalyticsMaturityReview,daSilvaCostaMamedeSilva2022RPAAdoptionSLR}
Die Priorisierung wird als Funktion aus Defizitstärke, Kontextrestriktionen und Abhängigkeiten modelliert, da Adoption- und Governance-Literatur sowohl organisatorische Voraussetzungen als auch Skalierungsmechanismen als begrenzende Faktoren beschreibt.\cite{TornatzkyFleischer1990TOE,KedzioraPenttinen2021RPAGovernance,daSilvaCostaMamedeSilva2022RPAAdoptionSLR}
Die Output-Spezifikation wird zudem so gestaltet, dass sie deterministische Felder und optionale textuelle Erläuterungen trennt, da Design-Science-Leitlinien eine klare Abgrenzung des Artefaktkerns von Darstellungs- oder Kommunikationskomponenten nahelegen.\cite{Hevner2004DesignScience,MarchSmith1995DesignNaturalScience}
Damit ist die Output-Struktur hinreichend formalisiert, um in Kapitel~4 als API- und Persistenzschema implementiert und in Kapitel~5 anhand definierter Qualitätskriterien evaluiert werden zu können.\cite{Hevner2004DesignScience,Nunamaker1990SystemsDevelopment}


\section{Konzeptionelles Informationsmodell}\label{sec:konzeptionellesInformationsmodell}

\subsubsection{Zweck und Abgrenzung}
Das konzeptionelle Informationsmodell präzisiert, welche Informationsobjekte das prototypische System benötigt, erzeugt und weiterverwendet, um die aus Nutzersicht relevanten Aufgaben systematisch zu unterstützen. \cite{normanDraper1986,iso9241_210_2019}
Damit fungiert das Informationsmodell als Brücke zwischen der Nutzungssicht und dem nachgelagerten Architekturentwurf, weil es die fachliche Struktur der Daten und Ergebnisse festlegt, auf die Komponentenverantwortlichkeiten und Kommunikationsbeziehungen später abgebildet werden. \cite{iso9241_210_2019,nunamaker1990}
Im Sinne des User Centered System Design wird diese Struktur nicht aus technischen Rahmenbedingungen abgeleitet, sondern aus den Zielen und Aufgaben der primären Nutzerrolle, da die Informationsbedarfe unmittelbar aus der Art der zu treffenden Entscheidungen und der dafür erforderlichen Nachvollziehbarkeit resultieren. \cite{normanDraper1986,iso9241_210_2019}

Das Modell wird als konzeptionelles UML-Klassendiagramm formuliert und beschreibt damit eine fachliche Sicht auf Struktur und Beziehungen der Informationsobjekte, ohne Persistenztechnologien, konkrete Datenbankschemata oder Implementierungsdetails festzulegen. \cite{omgUML251}
Auf dieser Grundlage kann das Architekturmodell im Anschluss Komponenten so schneiden, dass Erhebung, Bewertung, Synthese und Maßnahmenkatalogerzeugung auf klar definierten Informationsobjekten operieren und die Konsistenz zur Nutzungssicht gewahrt bleibt. \cite{nunamaker1990,iso9241_210_2019}


\subsubsection{UML-Klassendiagramm}
Abbildung X zeigt das konzeptionelle UML-Klassendiagramm als zentrale Darstellung des Informationsmodells. Das Diagramm ist in vier zusammenhängende Bereiche gegliedert, die den im System vorgesehenen Ablauf aus Nutzersicht abbilden. Der Block zur Erhebung umfasst die fachliche Beschreibung des Use Cases sowie die strukturierte Datenerhebung über einen Fragenkatalog mit Fragen und zugehörigen Antworten, die in einem Antwortsatz gebündelt werden. Darauf aufbauend beschreibt der Block zur Bewertung die getrennte Erzeugung zweier Ergebnisse in Form einer \gls{bi}-bezogenen Reifegradeinschätzung und einer Reifegradeinschätzung für Prozessautomatisierung. Der Block zur Synthese integriert beide Teilbewertungen in einer konsolidierten Ergebnisstruktur, die als Grundlage für die Maßnahmenableitung dient. Der Block zu den Maßnahmen repräsentiert den erzeugten Maßnahmenkatalog, die enthaltenen Maßnahmen sowie die nutzerseitige Finalisierung im Sinne der assistierenden Systemausprägung.

\begin{figure}[H]
\centering
\includegraphics[width=14cm]{arbeit/images/UML-Klassendiagramm.png}
\caption[Klassendiagramm]{UML-Klassendiagramm zur Darstellung der fachlichen Struktur der Informationsobjekte}
\label{IMG:UML-Klassendiagramm}
\end{figure}

Der zentrale Datenfluss beginnt mit der Erstellung eines Use Cases, der als fachlicher Ausgangspunkt für die Analyse dient. Auf dieser Grundlage wird ein Antwortsatz erzeugt, der die Antworten auf den zugeordneten Fragenkatalog enthält und damit die strukturierte Eingabebasis für die weitere Verarbeitung bildet. Der Antwortsatz wird anschließend für zwei getrennte Bewertungen verwendet, sodass sowohl eine \gls{bi}-Reifegradeinschätzung als auch eine Automatisierungsreifegradeinschätzung erzeugt werden können. Beide Bewertungen fließen in die Synthese ein, in der die Ergebnisse zusammengeführt und für die Ableitung von Handlungsempfehlungen vorbereitet werden. Aus der Synthese wird schließlich ein Maßnahmenkatalog abgeleitet, der konkrete Maßnahmen enthält und durch eine optionale Nutzerfinalisierung ergänzt wird, wodurch die Entscheidungshoheit bei der Nutzerrolle verbleibt.


\subsubsection{Erläuterung der zentralen Klassen und Beziehungen}
Die Klasse \texttt{UseCase} bildet den fachlichen Ausgangspunkt des Informationsmodells und repräsentiert den vom Nutzer beschriebenen Anwendungsfall, der im System analysiert werden soll. Über das Attribut \texttt{type} wird festgehalten, ob der Use Case einen \gls{bi}-Schwerpunkt, einen Automatisierungsschwerpunkt oder eine kombinierte Betrachtung besitzt. Diese Einordnung dient im weiteren Verlauf dazu, die Bewertung und Maßnahmenableitung konzeptionell konsistent an der ursprünglichen Zielsetzung des Nutzers auszurichten.

Die strukturierte Erhebung wird durch die Klassen \texttt{Questionnaire}, \texttt{Question}, \texttt{AnswerSet} und \texttt{Answer} abgebildet. \texttt{Questionnaire} repräsentiert den Fragenkatalog als konzeptionelles Erhebungsinstrument, während \texttt{Question} die einzelnen Fragen enthält, einschließlich der Kennzeichnung, ob eine Frage verpflichtend zu beantworten ist. \texttt{AnswerSet} bündelt die vom Nutzer gegebenen Antworten und fungiert damit als konsolidierte Eingabestruktur für die nachgelagerte Verarbeitung. Die Beziehung zwischen \texttt{Answer} und \texttt{Question} stellt sicher, dass jede Antwort eindeutig auf eine konkrete Frage referenziert und damit semantisch interpretierbar bleibt. Die Verknüpfung zwischen \texttt{UseCase} und \texttt{AnswerSet} drückt aus, dass genau ein Antwortsatz zu einem konkreten Use Case gehört und die Erhebung damit stets in den Kontext einer spezifischen Problemstellung eingebettet ist. Der Status des \texttt{AnswerSet} dient als konzeptioneller Mechanismus, um zwischen einem vorläufigen und einem validierten Zustand der Eingaben zu unterscheiden, ohne hierbei bereits technische Validierungsverfahren festzulegen.

Die Bewertung wird über die abstrakte Klasse \texttt{Assessment} und ihre Spezialisierungen \texttt{BIAssessment} und \texttt{PAAssessment} modelliert. \texttt{Assessment} bündelt die gemeinsamen Eigenschaften der Bewertungsergebnisse, insbesondere eine übergreifende Einordnung \texttt{overallLevel} sowie eine zusammenfassende Begründung \texttt{summary}. Die Spezialisierungen stehen für die getrennte Reifegradbewertung in den beiden Domänen und operationalisieren damit das in der Konzeption festgelegte Prinzip, dass BI und Prozessautomatisierung zunächst getrennt eingeschätzt werden, bevor eine integrierte Ableitung erfolgt. Die Beziehungen \texttt{AnswerSet} zu \texttt{BIAssessment} sowie \texttt{AnswerSet} zu \texttt{PAAssessment} legen fest, dass beide Bewertungen auf derselben Erhebungsbasis beruhen, wodurch Vergleichbarkeit und Konsistenz der Eingaben über beide Teilbewertungen hinweg gewährleistet werden.

Die Klasse \texttt{Synthesis} stellt die Integrationsebene dar und verbindet beide Teilbewertungen in einer konsolidierten Sicht. Die Beziehungen \texttt{Synthesis} zu \texttt{BIAssessment} und \texttt{Synthesis} zu \texttt{PAAssessment} verdeutlichen, dass die Synthese nicht als zusätzliche Bewertung neben den Teilbewertungen zu verstehen ist, sondern als Zusammenführung ihrer Ergebnisse. Damit wird konzeptionell abgesichert, dass die Ableitung von Maßnahmen nicht ausschließlich aus einer einzelnen Domänenperspektive erfolgt, sondern auf einer integrierten Interpretation beider Reifegrade beruht.

Die Ergebnisableitung wird durch \texttt{MeasureCatalog} und \texttt{Measure} repräsentiert. Der \texttt{MeasureCatalog} bündelt die vorgeschlagenen Maßnahmen als zentrales Ergebnisartefakt und trägt mit dem Attribut \texttt{status} den Zustand, ob es sich um einen Entwurf oder eine finalisierte Fassung handelt. Die Klasse \texttt{Measure} beschreibt einzelne Maßnahmen in einer für strategische Planung geeigneten Granularität, wobei \texttt{prioritySuggestion} eine vom System abgeleitete Priorisierung abbildet und \texttt{category} die Zuordnung zu BI, Automatisierung oder einem kombinierten Charakter ermöglicht. Die Beziehung \texttt{Synthesis} zu \texttt{MeasureCatalog} drückt aus, dass der Maßnahmenkatalog aus der integrierten Sicht abgeleitet wird und damit unmittelbar an die zuvor erzeugten Bewertungsergebnisse gekoppelt ist.

Die assistierende Systemausprägung wird schließlich über \texttt{UserSelection} abgebildet. Diese Klasse modelliert die nutzerseitige Finalisierung als optionale Ergänzung zu einer Maßnahme und hält fest, ob eine Maßnahme ausgewählt wird und welche finale Priorität der Nutzer vergibt. Die Kardinalität \texttt{Measure} 0..1 ↔ 0..1 \texttt{UserSelection} verdeutlicht, dass die Nutzerfinalisierung nicht zwingend erfolgen muss, aber als konzeptionelle Möglichkeit vorgesehen ist, um die Entscheidungshoheit beim Nutzer zu verankern. Dadurch entsteht eine klare Trennung zwischen einer systemseitig vorgeschlagenen Priorität und einer nutzerseitig bestätigten oder angepassten Priorisierung, ohne die interne Logik der Vorschlagserzeugung in diesem Modell vorwegzunehmen.


\subsubsection{Übergang zum konzeptionellen Architekturmodell}
Das konzeptionelle Informationsmodell legt die Informationsobjekte fest, die über den Verlauf einer Analyse hinweg vorgehalten werden müssen, um eine konsistente Bearbeitung und Nachvollziehbarkeit der Ergebnisse sicherzustellen. Persistiert werden in konzeptioneller Sicht insbesondere der \texttt{UseCase} als fachlicher Ausgangspunkt sowie der zugehörige \texttt{AnswerSet} mit den referenzierten Antworten, da diese Objekte die vollständige Eingabebasis und den Kontext der Bewertung abbilden. Ergänzend ist eine persistente Ablage des \texttt{MeasureCatalog} einschließlich der enthaltenen \texttt{Measure}-Objekte sowie einer optionalen \texttt{UserSelection} sinnvoll, weil der Maßnahmenkatalog als zentrales Ergebnisartefakt der Nutzung dient und eine spätere Weiterverwendung oder erneute Überarbeitung unterstützen soll.

Demgegenüber entstehen die Bewertungsergebnisse und die integrierte Sicht durch Verarbeitungsschritte auf Basis der erhobenen Eingaben. Die beiden \texttt{Assessment}-Ausprägungen \texttt{BIAssessment} und \texttt{PAAssessment} werden aus dem validierten Antwortsatz abgeleitet und bilden die domänenspezifischen Ergebnisse der Reifegradbewertung ab. Darauf aufbauend wird die \texttt{Synthesis} als zusammenführendes Ergebnis erzeugt, aus dem anschließend der \texttt{MeasureCatalog} abgeleitet wird. Diese Abfolge verdeutlicht, dass wesentliche Ergebnisobjekte nicht direkt durch den Nutzer erfasst werden, sondern als systemseitige Artefakte aus der Verarbeitung entstehen.

Das nachfolgende konzeptionelle Architekturmodell greift diese Trennung zwischen persistenten Basisobjekten und verarbeiteten Ergebnisobjekten auf und beschreibt, welche Komponenten für Erhebung, Bewertung, Synthese und Maßnahmenableitung verantwortlich sind. Dabei wird insbesondere festgelegt, wie die Verarbeitungsschritte auf Komponenten verteilt werden und welche Kommunikationsbeziehungen notwendig sind, um die im Informationsmodell definierten Objekte zu erzeugen, zu lesen und weiterzuverarbeiten.


\section{Konzeptionelles Architekturmodell}\label{sec:konzeptionellesArchitekturmodell}
\subsubsection{Zweck, Architektursicht und Abgrenzung}
Das konzeptionelle Architekturmodell ergänzt die Nutzungssicht und das Informationsmodell um eine strukturierende Systemperspektive, indem es beschreibt, wie die im Anwendungsfallmodell abgeleiteten Systemleistungen in Verantwortungsbereiche zerlegt und als Komponenten organisiert werden. Diese Architektursicht ist im Rahmen des Theory Building erforderlich, weil das Artefakt nicht nur über seine Ergebnisse, sondern auch über eine nachvollziehbare innere Struktur beschrieben werden muss, die die Erzeugung dieser Ergebnisse plausibel erklärt. \cite{nunamaker1990}
Während das Informationsmodell festlegt, welche Informationsobjekte im Systemverlauf entstehen und persistiert werden, adressiert die Architekturmodellierung die Frage, welche logischen Bausteine diese Objekte erzeugen, lesen und weiterverarbeiten und wie die Zusammenarbeit dieser Bausteine konzeptionell ausgestaltet ist.

In diesem Abschnitt wird daher eine konzeptionelle Ebene beschrieben, die zwischen fachlicher Anforderung und technischer Realisierung liegt. Der Fokus liegt auf einer logischen Komponentenzerlegung sowie auf der Kommunikation zwischen diesen Komponenten entlang eines typischen End-to-End-Ablaufs. Das Modell abstrahiert dabei bewusst von konkreten Produktentscheidungen und Implementierungstechnologien, um die Lösung unabhängig von spezifischen Tool- oder Plattformvorgaben zu formulieren.

Der prototypische Charakter des Systems bleibt als zentrale Prämisse bestehen. Das System wird als geschlossenes System ohne direkte Integrationen in operative Unternehmenssysteme modelliert, sodass sämtliche Eingaben über die Interaktion der primären Nutzerrolle erfolgen und Ausgaben innerhalb des Systems erzeugt und bereitgestellt werden. Externe Systeme werden daher im Architekturmodell nicht als Interaktionspartner abgebildet, sondern die Systemgrenze wird strikt auf die internen Komponenten gelegt.

Gegenstand der folgenden Modellierung sind ausschließlich das UML-Komponentenmodell sowie die UML-Darstellung der Kommunikationsbeziehungen zwischen den Komponenten. Nicht Bestandteil dieses Abschnitts sind Deployment- oder Infrastrukturentscheidungen, konkrete Ausprägungen des Betriebsmodells, Sicherheits- und Monitoring-Implementierungen oder detaillierte nichtfunktionale Auslegungen, sofern diese nicht unmittelbar für das Verständnis der konzeptionellen Struktur und der Kommunikationslogik erforderlich sind.


\subsubsection{Komponentenmodell als Struktur des Systems}
Abbildung X zeigt das konzeptionelle UML-Komponentenmodell des prototypischen Systems als logische Strukturierung der im Anwendungsfall- und Informationsmodell beschriebenen Systemleistungen. Die Darstellung ist als Pipeline entlang des zentralen Datenflusses aufgebaut, sodass die Komponenten unmittelbar den Verarbeitungsstufen \textit{Erhebung}, \textit{Bewertung}, \textit{Synthese} und \textit{Maßnahmenableitung} zugeordnet werden können. Zusätzlich sind Querschnittskomponenten ausgewiesen, die keine eigenständige fachliche Stufe bilden, aber die Ausführung der Pipeline unterstützen, insbesondere die konzeptionelle Datenhaltung sowie die assistierende Verarbeitung durch das \gls{llm}. Der Zweck der Modellierung liegt darin, Verantwortlichkeiten zu entkoppeln und die Erzeugung der im Informationsmodell definierten Objekte als Ergebnis wohldefinierter Komponenteninteraktionen erklärbar zu machen.

\begin{figure}[H]
\centering
\includegraphics[width=14cm]{arbeit/images/UML_Architekturmodell.png}
\caption[Komponentendiagramm]{UML-Komponentendiagramm zur Darstellung der konzeptionellen Architektur}
\label{IMG:UML-Architekturmodell}
\end{figure}

Die Komponente \textit{UI/Interaction} bildet die Schnittstelle zur primären Nutzerrolle und bündelt alle Interaktionsschritte, die in der Erhebungs- und Ergebnisphase erforderlich sind. In der Erhebung unterstützt sie die strukturierte Eingabe des \texttt{UseCase} und die Beantwortung des Fragenkatalogs, indem sie Nutzerführung und Eingabeerfassung bereitstellt. In der Ergebnisarbeit stellt sie Bewertungen, Synthese und Maßnahmenkatalog dar und ermöglicht optional die Finalisierung, indem eine nutzerseitige Auswahl beziehungsweise Priorisierung einzelner Maßnahmen als \texttt{UserSelection} erfasst wird. Die Komponente übernimmt damit die Verantwortung, Nutzeraktionen in fachliche Eingaben zu überführen und systemseitige Ergebnisse in einer für strategische Entscheidungen verwendbaren Form zurückzugeben, ohne selbst Bewertungslogik zu enthalten.

Der \textit{Questionnaire Service} ist für die Durchführung des strukturierten Fragenkatalogs verantwortlich und bildet die zentrale Komponente der Erhebungsstufe. Er verwaltet die Struktur des \texttt{Questionnaire} einschließlich der \texttt{Question}-Elemente und erzeugt aus den Nutzerantworten den konsolidierten \texttt{AnswerSet}. In dieser Komponente ist zudem die Validierung der Eingaben verortet, da die Unterscheidung zwischen einem vorläufigen und einem validierten Antwortsatz im Informationsmodell explizit vorgesehen ist. Der Service stellt damit sicher, dass die nachgelagerten Bewertungskomponenten eine stabile Eingabebasis erhalten, ohne dass bereits eine spezifische Implementierung der Validierungslogik festgelegt wird.

Die Bewertungsstufe ist in zwei getrennte Komponenten aufgeteilt, um die konzeptionell gewählte Trennung zwischen \gls{bi}- und Automatisierungsreifegradbewertung abzubilden. Der \textit{Assessment Service BI} erzeugt auf Basis eines validierten \texttt{AnswerSet} ein \texttt{BIAssessment} und verantwortet damit die domänenspezifische Interpretation der Eingaben im Kontext der BI-Fähigkeiten. Der \textit{Assessment Service PA} erzeugt analog ein \texttt{PAAssessment} und fokussiert die Reifegradbewertung im Kontext der Prozessautomatisierung. Die Trennung in zwei Komponenten reduziert Kopplung, erleichtert eine unabhängige Weiterentwicklung der Bewertungslogiken und stellt sicher, dass beide Teilbewertungen als eigenständige Ergebnisartefakte für die nachgelagerte Synthese verfügbar bleiben.

Der \textit{Synthesis Service} bildet die Integrationsstufe der Pipeline. Er nimmt beide Teilbewertungen entgegen und erzeugt daraus ein \texttt{Synthesis}-Objekt, das die Ergebnisse in einer konsolidierten Sicht zusammenführt. Inhaltlich liegt die Verantwortung dieser Komponente darin, Abhängigkeiten und Wechselwirkungen zwischen BI- und Automatisierungsperspektive so aufzubereiten, dass daraus eine kohärente Grundlage für die Maßnahmenableitung entsteht. Der Service ist damit bewusst von den domänenspezifischen Bewertungen getrennt, um die Synthese als eigenständige Systemleistung und nicht als Nebenprodukt einer Einzelbewertung zu modellieren.

Aufbauend auf der Synthese erzeugt der \textit{Recommendation Service} den \texttt{MeasureCatalog} als zentrales Ergebnisartefakt und enthält die Logik zur Ableitung und Strukturierung konkreter Maßnahmen. Die Komponente verantwortet die initiale Priorisierung als Vorschlag und berücksichtigt, dass der Maßnahmenkatalog in der assistierenden Systemausprägung durch den Nutzer finalisiert werden kann. Entsprechend ist hier auch die Verarbeitung nutzerseitiger Anpassungen verortet, sodass die finale Fassung des Katalogs als konsistentes Ergebnisobjekt vorliegt.

Als Querschnittskomponente ist das \textit{Repository} modelliert, das die konzeptionelle Datenhaltung für die wesentlichen Informationsobjekte bereitstellt. Dazu gehören insbesondere \texttt{UseCase} und \texttt{AnswerSet} als Eingabebasis sowie die durch Verarbeitung erzeugten Ergebnisobjekte \texttt{BIAssessment}, \texttt{PAAssessment}, \texttt{Synthesis} und \texttt{MeasureCatalog} einschließlich optionaler \texttt{UserSelection}. Die Modellierung als eigene Komponente dient der klaren Trennung zwischen fachlicher Verarbeitung und Datenhaltung, ohne konkrete Persistenztechnologien festzulegen.

Der \textit{LLM Service} ist als interne, assistierende Komponente ausgewiesen, die von mehreren Stufen der Pipeline genutzt werden kann. Er übernimmt keine eigenständige fachliche Verantwortung im Sinne eines separaten Prozessschritts, sondern unterstützt insbesondere bei der semantischen Interpretation von Texteingaben sowie bei der Formulierung von Begründungen und Maßnahmenvorschlägen. Die Platzierung als Querschnitt unterstreicht, dass das \gls{llm} innerhalb der Systemgrenze agiert und die Entscheidungshoheit beim Nutzer verbleibt, während die fachliche Orchestrierung weiterhin durch die domänenspezifischen Services erfolgt.


\subsubsection{Kommunikation zwischen Komponenten}
Die Kommunikation zwischen den Komponenten wird in diesem Abschnitt entlang des im Komponentenmodell dargestellten Hauptflusses erläutert. Abbildung X strukturiert die Interaktionen als Pipeline von der Erhebung über die domänenspezifischen Bewertungen und die Synthese bis zur Ableitung des Maßnahmenkatalogs. Der dargestellte Nachrichtenfluss ist dabei als Referenzinteraktion zu verstehen, die den typischen End-to-End-Durchlauf einer Analyse abbildet und unmittelbar an die im Informationsmodell definierten Objekte anknüpft.

Der Ablauf beginnt in der Komponente \textit{UI/Interaction}, in der der Nutzer den zu analysierenden \texttt{UseCase} erfasst. Diese Informationen werden an den \textit{Questionnaire Service} übergeben, der die Erhebung strukturiert und die Antworten des Nutzers in einem \texttt{AnswerSet} konsolidiert. Der \textit{Questionnaire Service} übernimmt zugleich die konzeptionelle Validierung der Eingaben, sodass der Antwortsatz als verlässliche Grundlage für die nachgelagerte Verarbeitung vorliegt. Damit entsteht eine klare Trennung zwischen der nutzerzentrierten Interaktionslogik in der UI und der fachlichen Erhebungs- und Prüflogik im Service.

Auf Basis des validierten \texttt{AnswerSet} werden anschließend zwei getrennte Bewertungsschritte angestoßen. Der \textit{Assessment BI} verarbeitet den Antwortsatz, um ein \texttt{BIAssessment} zu erzeugen, während der \textit{Assessment PA} parallel oder sequenziell ein \texttt{PAAssessment} ableitet. Diese Aufteilung entspricht der konzeptionellen Vorgabe, BI und Prozessautomatisierung zunächst domänenspezifisch zu bewerten, bevor eine integrierte Betrachtung erfolgt. Beide Bewertungsergebnisse werden im nächsten Schritt an den \textit{Synthesis Service} übergeben, der daraus ein \texttt{Synthesis}-Objekt erzeugt und damit die beiden Teilbewertungen in einer konsolidierten Sicht zusammenführt. Diese Kommunikation bildet den Übergang von domänenspezifischen Ergebnisartefakten zu einer integrierten Grundlage für die Maßnahmenableitung ab.

Der \textit{Recommendation Service} nutzt die erzeugte \texttt{Synthesis}, um den \texttt{MeasureCatalog} abzuleiten und in eine für den Nutzer interpretierbare Struktur zu überführen. Der Maßnahmenkatalog wird anschließend an die \textit{UI/Interaction} zurückgegeben, die das Ergebnis darstellt und eine Weiterbearbeitung ermöglicht. Im Sinne der assistierenden Systemausprägung ist eine optionale Nutzeranpassung vorgesehen, die sich in einer \texttt{UserSelection} ausdrückt. Diese Nutzerfinalisierung ergänzt den Maßnahmenkatalog um eine explizite Auswahl oder Priorisierung durch den Nutzer, ohne die systemseitig generierte Vorschlagslogik zu ersetzen.

Die im Komponentenmodell als Querschnitt ausgewiesenen Elemente wirken unterstützend, ohne den Hauptfluss zu erweitern. Das \textit{Repository} dient konzeptionell der Vorhaltung sowohl der Eingabeobjekte (\texttt{UseCase}, \texttt{AnswerSet}) als auch der erzeugten Ergebnisobjekte (\texttt{BIAssessment}, \texttt{PAAssessment}, \texttt{Synthesis}, \texttt{MeasureCatalog}) einschließlich der optionalen \texttt{UserSelection}. Damit ist sichergestellt, dass Ergebnisse nachvollziehbar bleiben und in späteren Sitzungen erneut genutzt oder angepasst werden können, ohne eine externe Systemintegration vorauszusetzen. Der \textit{LLM Service} ist als interne Assistenzkomponente modelliert, die insbesondere bei der semantischen Interpretation von Eingaben sowie bei der Formulierung von Begründungen und Maßnahmen unterstützt. Seine Rolle bleibt dabei auf Unterstützung beschränkt, da die fachliche Verantwortung weiterhin bei den jeweiligen Verarbeitungskomponenten liegt und die finale Entscheidung über Maßnahmen bei der Nutzerrolle verbleibt.


\subsubsection{Architekturentscheidungen und Zuordnung zu Informationsobjekten}
Die zentrale Architekturentscheidung besteht in einer komponentenbasierten Zerlegung entlang der fachlichen Verarbeitungsstufen, um Verantwortlichkeiten eindeutig zuzuordnen und die aus der Nutzungssicht abgeleiteten Systemleistungen in voneinander abgrenzbaren Bausteinen abzubilden. Die \textit{UI/Interaction} wird dabei strikt auf Interaktion und Ergebnisdarstellung begrenzt und enthält keine domänenspezifische Bewertungslogik. Diese Trennung unterstützt die Nachvollziehbarkeit des Systems, da fachliche Verarbeitungsschritte in dedizierten Services verankert werden und die UI als reine Vermittlungsschicht zwischen Nutzer und System fungiert.

Der \textit{Questionnaire Service} bildet den Verantwortungsbereich der Erhebung und Validierung. Ihm sind die Informationsobjekte \texttt{UseCase}, \texttt{Questionnaire}, \texttt{Question}, \texttt{AnswerSet} und \texttt{Answer} zugeordnet, wobei insbesondere die Konsolidierung und der Validierungsstatus des \texttt{AnswerSet} als Übergabepunkt an die Bewertungskomponenten dienen. Die Bewertung wird bewusst in \textit{Assessment BI} und \textit{Assessment PA} getrennt, um die konzeptionell geforderte domänenspezifische Betrachtung abzubilden und eine Vermischung von Bewertungslogiken zu vermeiden. Entsprechend erzeugt \textit{Assessment BI} das Objekt \texttt{BIAssessment}, während \textit{Assessment PA} das Objekt \texttt{PAAssessment} erzeugt, jeweils ausschließlich auf Basis des validierten \texttt{AnswerSet}.

Die \textit{Synthesis Service} wird als eigener Verantwortungsbereich modelliert, um die Integration der Teilbewertungen als eigenständige Systemleistung zu verankern. Dieser Service erzeugt das Objekt \texttt{Synthesis} und nutzt dazu beide Assessments, ohne deren domänenspezifische Ergebnisse zu verändern. Auf dieser integrierten Grundlage übernimmt der \textit{Recommendation Service} als separater Verantwortungsbereich die Ableitung der Handlungsempfehlungen. Ihm sind \texttt{MeasureCatalog} und \texttt{Measure} zugeordnet; die optionale Nutzerfinalisierung wird über \texttt{UserSelection} abgebildet, wodurch die assistierende Systemausprägung explizit im Modell verankert bleibt.

Das \textit{Repository} wird als Querschnittskomponente geführt, weil sowohl Eingabe- als auch Ergebnisobjekte über den End-to-End-Ablauf hinweg verfügbar sein müssen. Konzeptionell werden dort \texttt{UseCase} und \texttt{AnswerSet} als Grundlage sowie \texttt{BIAssessment}, \texttt{PAAssessment}, \texttt{Synthesis}, \texttt{MeasureCatalog} und optional \texttt{UserSelection} als erzeugte Ergebnisartefakte vorgehalten. Der \textit{LLM Service} wird ebenfalls als Querschnitt verstanden, da er mehrere Verarbeitungsschritte unterstützen kann, ohne als eigenständige fachliche Stufe aufzutreten und ohne die Verantwortungszuordnung der fachlichen Services zu ersetzen.



\section*{Zusammenfassung}
Kapitel 3 beschreibt das prototypische System als konzeptionelles Artefakt in drei aufeinander aufbauenden Sichten und stellt damit eine konsistente Grundlage für die weitere Ausgestaltung und Bewertung bereit. Ausgangspunkt ist eine nutzerzentrierte Perspektive, in der die Aufgaben und Entscheidungsbedarfe der primären Nutzerrolle die Ableitung der Systemleistungen leiten. Das konzeptionelle Anwendungsfallmodell formalisiert die Systemgrenze und den End-to-End-Ablauf von der Use-Case-Erfassung und der strukturierten Befragung bis zur Generierung eines Maßnahmenkatalogs, wobei ein \gls{llm} assistierend eingesetzt wird und die Entscheidungshoheit beim Nutzer verbleibt. Das konzeptionelle Informationsmodell präzisiert die hierfür erforderlichen Informationsobjekte und bildet den Datenfluss von \texttt{UseCase} und \texttt{AnswerSet} über getrennte BI- und Automatisierungsbewertungen bis zur Synthese und zum Maßnahmenkatalog ab. Das konzeptionelle Architekturmodell überführt diese Logik in eine Komponentenpipeline mit klaren Verantwortlichkeiten und einer nachvollziehbaren Kommunikation entlang der Referenzinteraktion.